\documentclass[11pt, letterpaper]{article}

% --- Idioma y codificación ---
\usepackage[spanish]{babel}
\usepackage[utf8]{inputenc}
\usepackage[T1]{fontenc}
\usepackage{lmodern}

% --- Matemáticas ---
\usepackage{amsmath, amssymb, amsthm}

% --- Formato ---
\usepackage[margin=2.5cm]{geometry}
\usepackage{parskip}
\usepackage{booktabs}
\usepackage{listings}
\usepackage{hyperref}
\hypersetup{colorlinks=true, linkcolor=blue!70!black, urlcolor=blue!70!black}

% --- Código ---
\lstset{
  language=C, basicstyle=\ttfamily\small,
  keywordstyle=\bfseries, commentstyle=\itshape,
  numbers=left, numberstyle=\tiny, frame=single,
  breaklines=true
}

% --- Entornos de teoremas ---
\theoremstyle{definition}
\newtheorem{definition}{Definición}[section]
\theoremstyle{plain}
\newtheorem{theorem}{Teorema}[section]
\newtheorem{lemma}[theorem]{Lema}
\newtheorem{corollary}[theorem]{Corolario}
\theoremstyle{remark}
\newtheorem{remark}{Observación}[section]

\title{%
  T01 --- Demostración formal:\\[4pt]
  Correctitud del heapsort
}
\author{Curso Linux--C--Rust --- B15 Estructuras de datos en C}
\date{}

\begin{document}
\maketitle
\tableofcontents

% ===================================================================
\section{El algoritmo}
% ===================================================================

Referencia del código cuya correctitud probamos:

\begin{lstlisting}
void heapsort(int arr[], int n) {
    // Fase 1: build max-heap
    for (int i = n / 2 - 1; i >= 0; i--)        // bucle B
        sift_down(arr, n, i);

    // Fase 2: extraer maximo
    for (int i = n - 1; i > 0; i--) {            // bucle E
        swap(&arr[0], &arr[i]);
        sift_down(arr, i, 0);
    }
}
\end{lstlisting}

La prueba se divide en tres partes: correctitud de la Fase~1 (build),
correctitud de la Fase~2 (extract), y complejidad.

% ===================================================================
\section{Definiciones}
% ===================================================================

\begin{definition}[Max-heap en subrango]
Decimos que \texttt{arr[0..m)} es un max-heap si para todo
$0 \leq j < m$:
\begin{itemize}
  \item Si $2j+1 < m$: $\texttt{arr}[j] \geq \texttt{arr}[2j+1]$
    \quad (padre $\geq$ hijo izquierdo).
  \item Si $2j+2 < m$: $\texttt{arr}[j] \geq \texttt{arr}[2j+2]$
    \quad (padre $\geq$ hijo derecho).
\end{itemize}
\end{definition}

\begin{definition}[Partición heap-sorted]
Decimos que el array tiene una partición~$(i)$ si:
\begin{enumerate}
  \item[\textbf{P1.}] \texttt{arr[0..i+1)} es un max-heap de tamaño
    $i+1$.
  \item[\textbf{P2.}] \texttt{arr[i+1..n)} contiene los $n-(i+1)$
    mayores elementos del array original, en orden ascendente.
  \item[\textbf{P3.}] Todo elemento en \texttt{arr[0..i+1)} es $\leq$
    todo elemento en \texttt{arr[i+1..n)}.
\end{enumerate}
\end{definition}

% ===================================================================
\section{Demostración 1: correctitud de la Fase 1 (build-heap)}
% ===================================================================

\begin{theorem}\label{thm:build}
Al terminar el bucle~B, \texttt{arr[0..n)} es un max-heap válido.
\end{theorem}

\begin{definition}[Invariante del bucle B]
\[
  \textbf{INV\_B}(i): \text{Para todo } j \text{ con } i < j < n,
  \text{ el nodo } j \text{ es raíz de un max-heap válido dentro de
  \texttt{arr[0..n)}}.
\]
\end{definition}

\begin{proof}[Prueba del Teorema~\ref{thm:build}]
Probamos el invariante en tres etapas.

\textbf{Inicialización} ($i = \lfloor n/2 \rfloor - 1$, antes de la
primera iteración). Los nodos con índice $j \geq \lfloor n/2 \rfloor$
son hojas. Una hoja es trivialmente un max-heap de tamaño~1 (no tiene
hijos que violen la propiedad). Para todo
$j > \lfloor n/2 \rfloor - 1$: el nodo~$j$ es hoja $\to$ raíz de un
max-heap válido. \checkmark

\textbf{Mantenimiento} (iteración~$i$). Al inicio de la iteración,
$\text{INV\_B}(i)$ dice que todos los nodos con índice $> i$ son
raíces de max-heaps válidos. En particular, los hijos de~$i$ (índices
$2i+1$ y $2i+2$, si existen) son raíces de max-heaps válidos.

Se ejecuta \texttt{sift\_down(arr, n, i)}. La operación sift-down
tiene la precondición de que ambos subárboles del nodo son max-heaps,
que se cumple por el invariante. Después de sift-down, el subárbol
enraizado en~$i$ es un max-heap válido.

Al decrementar $i$ a $i-1$, el invariante $\text{INV\_B}(i-1)$ dice
que todos los nodos con índice $> i - 1$ (es decir, $\geq i$) son
raíces de max-heaps. Como acabamos de reparar~$i$, esto se cumple.
\checkmark

\textbf{Terminación} ($i = -1$). El bucle termina cuando $i < 0$. El
invariante $\text{INV\_B}(-1)$ dice que todos los nodos con índice
$> -1$ (es decir, todos los nodos $0, 1, \ldots, n{-}1$) son raíces
de max-heaps válidos. En particular, el nodo~0 (la raíz) es raíz de
un max-heap válido sobre \texttt{arr[0..n)}.
\end{proof}

% ===================================================================
\section{Demostración 2: correctitud de la Fase 2 (extract)}
% ===================================================================

\begin{theorem}\label{thm:extract}
Al terminar el bucle~E, \texttt{arr[0..n)} está ordenado en orden
ascendente.
\end{theorem}

\begin{definition}[Invariante del bucle E]
$\textbf{INV\_E}(i)$: al inicio de cada iteración con valor~$i$:
\begin{enumerate}
  \item[\textbf{E1.}] \texttt{arr[0..i+1)} es un max-heap de tamaño
    $i+1$.
  \item[\textbf{E2.}] \texttt{arr[i+1..n)} contiene los $n-(i+1)$
    mayores elementos del array original, en orden ascendente.
  \item[\textbf{E3.}] Todo elemento en \texttt{arr[0..i+1)} es $\leq$
    todo elemento en \texttt{arr[i+1..n)}.
\end{enumerate}
\end{definition}

\begin{proof}[Prueba del Teorema~\ref{thm:extract}]
\textbf{Inicialización} ($i = n - 1$, antes de la primera iteración).
\begin{itemize}
  \item (E1): \texttt{arr[0..n)} es un max-heap de tamaño~$n$ (por la
    Fase~1). \checkmark
  \item (E2): \texttt{arr[n..n)} es vacío. La condición se satisface
    trivialmente. \checkmark
  \item (E3): No hay elementos en \texttt{arr[n..n)}, condición
    trivial. \checkmark
\end{itemize}

\textbf{Mantenimiento} (iteración~$i$). Asumimos $\text{INV\_E}(i)$.

\emph{Paso~1: \texttt{swap(\&arr[0], \&arr[i])}.} Por~(E1),
\texttt{arr[0]} es el máximo de \texttt{arr[0..i+1)} (propiedad del
max-heap: la raíz es el mayor). Por~(E3), este máximo es $\leq$ todos
los elementos en \texttt{arr[i+1..n)}. Después del swap:
\texttt{arr[i]} contiene el antiguo máximo de \texttt{arr[0..i+1)}, y
\texttt{arr[0]} contiene el antiguo \texttt{arr[i]}.

\emph{Paso~2: \texttt{sift\_down(arr, i, 0)}.} Después del swap,
\texttt{arr[0..i)} puede no ser un max-heap (porque \texttt{arr[0]}
fue reemplazado). Sin embargo, los subárboles de la raíz (enraizados
en índices~1 y~2) siguen siendo max-heaps válidos: el swap solo afectó
a \texttt{arr[0]} y \texttt{arr[i]}, y el nodo \texttt{arr[i]} ya no
es parte de la zona de heap \texttt{[0..i)}. La precondición de
sift-down se cumple. Después de \texttt{sift\_down(arr, i, 0)},
\texttt{arr[0..i)} es un max-heap de tamaño~$i$.

\emph{Verificación del invariante para $i - 1$:}
\begin{itemize}
  \item (E1): \texttt{arr[0..i)} = \texttt{arr[0..(i{-}1)+1)} es un
    max-heap de tamaño $i = (i{-}1)+1$. \checkmark
  \item (E2): \texttt{arr[i..n)} contiene el antiguo máximo de
    \texttt{arr[0..i+1)} (ahora en \texttt{arr[i]}) seguido de
    \texttt{arr[i+1..n)} (que por hipótesis contenía los $n-i-1$
    mayores en orden). El antiguo máximo es, por~(E3), $\leq$ todos
    los elementos en \texttt{arr[i+1..n)}, y es $\geq$ todos los de
    \texttt{arr[0..i+1)}. Así, $\texttt{arr}[i] \leq \texttt{arr}[i+1]$,
    y \texttt{arr[i..n)} contiene los $n-i$ mayores en orden
    ascendente. \checkmark
  \item (E3): El máximo de \texttt{arr[0..i)} (que es \texttt{arr[0]}
    tras sift-down) era un elemento de \texttt{arr[0..i+1)} antes del
    swap. Por~(E3) del paso anterior, era $\leq$ todos los de
    \texttt{arr[i+1..n)}. Como \texttt{arr[i]} es el antiguo máximo de
    \texttt{arr[0..i+1)}, tenemos
    $\texttt{arr}[0] \leq \texttt{arr}[i] \leq \texttt{arr}[i+1]$.
    \checkmark
\end{itemize}

\textbf{Terminación} ($i = 0$). El bucle termina cuando $i = 0$.
El invariante $\text{INV\_E}(0)$ dice:
\begin{itemize}
  \item (E1): \texttt{arr[0..1)} es un max-heap de tamaño~1
    (trivialmente cierto).
  \item (E2): \texttt{arr[1..n)} contiene los $n-1$ mayores elementos
    en orden ascendente.
  \item (E3): $\texttt{arr}[0] \leq$ todos los elementos en
    \texttt{arr[1..n)}.
\end{itemize}

Combinando: \texttt{arr[0]} es el menor elemento, y \texttt{arr[1..n)}
son los $n-1$ mayores en orden ascendente. Por lo tanto,
\texttt{arr[0..n)} está completamente ordenado.
\end{proof}

% ===================================================================
\section{Demostración 3: complejidad
  \texorpdfstring{$\Theta(n \log n)$}{Theta(n log n)}}
% ===================================================================

\begin{theorem}\label{thm:complexity}
Heapsort tiene complejidad $\Theta(n \log n)$ en el peor caso.
\end{theorem}

\begin{proof}
\emph{Cota superior.} La Fase~1 (build-heap) cuesta $O(n)$ (demostrado
en T04/README\_EXTRA). La Fase~2 ejecuta $n-1$ extracciones, cada una
con un swap $O(1)$ y un sift-down $O(\log i)$ (la zona de heap tiene
tamaño~$i$):
\[
  T_2 = \sum_{i=1}^{n-1} O(\log i)
  \leq \sum_{i=1}^{n-1} O(\log n) = O(n \log n)
\]
Total: $O(n) + O(n \log n) = O(n \log n)$.

\emph{Cota inferior.} Heapsort es un algoritmo de ordenamiento basado
en comparaciones. Por la cota inferior de $\Omega(n \log n)$ para
ordenamiento por comparación (demostrada en C08/README\_EXTRA), todo
algoritmo de este tipo requiere $\Omega(n \log n)$ comparaciones en el
peor caso.

Combinando: heapsort es $\Theta(n \log n)$ en el peor caso.
\end{proof}

\begin{corollary}[Complejidad en todos los casos]
\mbox{}
\end{corollary}

\begin{table}[h]
\centering
\begin{tabular}{lcp{6cm}}
\toprule
\textbf{Caso} & \textbf{Complejidad} & \textbf{Razón} \\
\midrule
Mejor    & $O(n \log n)$
  & Incluso si el array ya está ordenado, la Fase~2 hace $n{-}1$
    sift-downs \\
Promedio & $O(n \log n)$
  & Los sift-downs en la Fase~2 descienden $\Theta(\log n)$ niveles
    en promedio \\
Peor     & $O(n \log n)$
  & Cada sift-down desciende a lo sumo
    $\lfloor \log_2 n \rfloor$ niveles \\
\bottomrule
\end{tabular}
\end{table}

\begin{remark}
Heapsort es uno de los pocos algoritmos con la \emph{misma}
complejidad asintótica en los tres casos.
\end{remark}

% ===================================================================
\section{Demostración 4: heapsort es in-place}
% ===================================================================

\begin{theorem}\label{thm:inplace}
Heapsort usa $O(1)$ memoria auxiliar.
\end{theorem}

\begin{proof}
El algoritmo opera sobre el array de entrada sin alocar estructuras
adicionales. Las únicas variables auxiliares son:
\begin{itemize}
  \item Índices del bucle ($i$, \texttt{index}, \texttt{left},
    \texttt{right}, \texttt{largest}): $O(1)$.
  \item Variable temporal para swap: $O(1)$.
\end{itemize}

No se usan arrays auxiliares ni recursión (sift-down es iterativo). La
versión recursiva de sift-down usaría $O(\log n)$ stack, pero la
versión iterativa (con \texttt{while}) es $O(1)$.
\end{proof}

% ===================================================================
\section{Resumen de resultados}
% ===================================================================

\begin{table}[h]
\centering
\begin{tabular}{llp{5.5cm}}
\toprule
\textbf{Resultado} & \textbf{Técnica} & \textbf{Clave} \\
\midrule
Fase~1 produce max-heap
  & Invariante de bucle~B
  & Hojas son heaps triviales, sift-down repara subárboles \\
Fase~2 produce array ordenado
  & Invariante de bucle~E (3 cláusulas)
  & Máximo del heap va a la zona ordenada \\
Complejidad $\Theta(n \log n)$
  & Suma + cota inferior
  & $O(n) + O(n\log n)$; $\Omega(n\log n)$ por comparaciones \\
In-place $O(1)$
  & Análisis de variables
  & Solo índices y temporal de swap \\
\bottomrule
\end{tabular}
\end{table}

\end{document}
